%% =====================================================================
%%  core/config/colores.tex  --  Paleta institucional
%% ---------------------------------------------------------------------
%%  Fuente de verdad: identidad-visual/tokens/colors.json (repo del sistema de diseño)
%%  (escalas puka / q'illu / neutral). Aqui solo se REPLICAN los valores;
%%  si el manual de marca cambia, actualizar alli primero y reflejar aqui.
%%
%%  Regla dura de la marca: SIN NARANJA. Nunca mezclar puka con qillu
%%  en degradados.
%% =====================================================================

\RequirePackage{xcolor}

% --- Escala puka (rojo cochinilla, ancla de marca = 600) -------------
\definecolor{puka600}{HTML}{A6122E}   % marca primaria
\definecolor{puka700}{HTML}{8A0F27}
\definecolor{puka800}{HTML}{6E0F22}   % "Vino"
\definecolor{puka900}{HTML}{4E0A18}

% --- Guinda del logo oficial de la Camara de Diputados ---------------
\definecolor{guindaoficial}{HTML}{4D1018}

% --- Escala q'illu (amarillo maiz, ancla = 400) ----------------------
\definecolor{qillu400}{HTML}{EAB61E}  % marca secundaria
\definecolor{qillu600}{HTML}{A67B10}

% --- Neutros calidos --------------------------------------------------
\definecolor{papel}{HTML}{F5F1EC}     % neutral.100
\definecolor{neutral500}{HTML}{A99E92}
\definecolor{neutral700}{HTML}{6E655D}
\definecolor{tinta}{HTML}{241F1C}     % neutral.900, texto principal
\definecolor{negro}{HTML}{000000}     % negro pleno (caja y filetes)

% --- Roles (que color usa cada cosa; cambiar aqui, no en los usos) ---
\colorlet{colormembrete}{negro}           % caja del nombre en el membrete
\colorlet{colorfilete}{negro}             % filetes horizontales (membrete y pie)
\colorlet{colortitulos}{guindaoficial}    % titulos de seccion (informes)
\colorlet{colorenlaces}{puka800}          % hyperref
\colorlet{colorsello}{guindaoficial}      % sello de consigna (punos + frase)

% --- Roles de la banda de prensa (comunicados/pronunciamientos) ------
%%  La franja "pallay" evoca el tejido ayacuchano: campo puka con
%%  rombos alternados. Solidos adyacentes (nunca degradados puka-qillu).
\colorlet{colorpallay}{puka600}           % campo de la franja textil
\colorlet{colorpallayacento}{qillu400}    % rombos de acento (alternan con blanco)
\colorlet{colorbandalinea}{qillu400}      % linea fina sobre la banda blanca
\colorlet{colorbandatag}{guindaoficial}   % cajita "DIPUTADO POR AYACUCHO"

% --- Escalas ampliadas (tokens/colors.json, pasos 50-300 y neutros) ---
%%  Agregadas 2026-09-03 para el INFORME DE SEMANA DE REPRESENTACION
%%  (classes/informe-sr.cls): fondos de cajas y fichas, filas alternas,
%%  bordes y chips. Regla de la marca: los pasos 50-300 JAMAS como texto.
\definecolor{puka50}{HTML}{FBEDEF}
\definecolor{puka100}{HTML}{F6D8DC}
\definecolor{puka200}{HTML}{EFB0B9}
\definecolor{qillu50}{HTML}{FDF6E3}
\definecolor{qillu100}{HTML}{FAECC4}
\definecolor{qillu200}{HTML}{F5DE97}
\definecolor{qillu500}{HTML}{C99614}
\definecolor{neutral200}{HTML}{E9E2D9}
\definecolor{neutral300}{HTML}{DCD2C6}
\definecolor{neutral400}{HTML}{C9BFB4}
\definecolor{neutral600}{HTML}{8A8079}
\definecolor{neutral800}{HTML}{524B45}
\definecolor{datoazul}{HTML}{4A5A66}   % azul-pizarra DATA-ONLY (6.a categoria de graficos)
\definecolor{exito}{HTML}{2E7D32}      % funcional (estados), nunca de marca

% --- Roles del informe de Semana de Representacion (informe-sr) ------
%%  El informe consume SOLO estos roles; cambiar aqui, no en la
%%  estructura. Modo "Blanco oficial" en el cuerpo; "Vino" en portada y
%%  separadores; Q'illu unicamente como filete/acento.
\colorlet{colorisrprimario}{puka600}      % kickers, numerales, acentos fuertes
\colorlet{colorisrprofundo}{puka800}      % portada, separadores, cajas oscuras (Vino)
\colorlet{colorisracento}{qillu400}       % filetes finos y acentos
\colorlet{colorisrtexto}{tinta}
\colorlet{colorisrtextosec}{neutral700}
\colorlet{colorisrfondo}{papel}           % fondo de fichas y cajas
\colorlet{colorisrfondosuave}{puka50}     % destacado suave (alertas, compromisos)
\colorlet{colorisrborde}{neutral300}
\colorlet{colorisrfila}{neutral200}       % filas alternas de tablas
